\chapter{Grundlagen}

\targets{
  \item Erstes Ziel
  \item Zweites Ziel
  \item Drittes Ziel
  \item Viertes Ziel
}

\website

\section{Was ist \LaTeX?}

\subsection{Einordnung}

\begin{Frame}{Dimensionen eines Dokumentes}
  \begin{description}
    \item[Inhalt] ist die \alert{Bedeutung} eines Textes
    \item[Struktur] ist der \alert{Aufbau} eines Textes
    \item[Form] ist das \alert{Aussehen} eines Textes
  \end{description}
\end{Frame}

\begin{Frame}{Struktur vs. Form}
  \begin{Beispiel}[Strukturelemente]
    \begin{itemize}
      \item Überschrift
      \item Listeneintrag
      \item Tabellenzelle
    \end{itemize}
  \end{Beispiel}

  \xxx

  \begin{Beispiel}[Formen]
    \begin{itemize}
      \item 16~pt, Arial, Fett, 2~em Abstand
      \item 2~cm Einzug, Bullet-Zeichen \textbullet\ am Zeilenanfang
      \item 3~cm breiter umrandeter Kasten
    \end{itemize}
  \end{Beispiel}
\end{Frame}

\begin{Frame}[fragile]{Seitenbeschreibungssprachen}
  \newcommand{\entry}[2]{\draw[maincolor, thick] (-.2,#1) -- (.2,#1) node[right] {\color{black}#2};}
  \hskip 8ex\begin{tikzpicture}
    \draw[maincolor, thick] (0,0) node[below] {\textbf{Form}} edge[<->] (0,6);
    \draw[maincolor] (0,6) node[above] {\textbf{Struktur}};
    \entry{.5}{Pixelgrafiken bzw. Paint};
    \entry{1}{Vektorgrafiken bzw. Inkscape};
    \entry{1.5}{PDF};
    \entry{2.5}{\TeX};
    \entry{3}{DTP-Tools wie z.\,B. Scribus};
    \entry{3.5}{Office Word, Writer, \ldots};
    \entry{4}{\alert{\LaTeX}};
    \entry{4.75}{HTML};
    \entry{5.5}{Outliner};
  \end{tikzpicture}
\end{Frame}

\begin{Frame}{\LaTeX}
  \begin{itemize}
    \item basiert auf dem Satzsystem \TeX
      \begin{itemize}
        \item \TeX\ wurde 1977 von Donald E. Knuth enwickelt
        \item Aktuelle Version: 3.1415926 (März 2008)
      \end{itemize}
    \item wurde Anfang der 1980er von Leslie Lamport entwickelt
    \item Aktuelle Version: 2011/06/27
  \end{itemize}
\end{Frame}

\subsection{Aufbau}

\begin{Frame}{Ein \TeX-Dokument}{Quelltext \texttt{story.tex}}
  \lstinputlisting[firstline=7]{demo/story.tex}
\end{Frame}

\begin{Frame}[fragile]{Ein \TeX-Dokument}{Kompilieren}
  \begin{lstlisting}[language={},morekeywords={tex,dvips,pstopdf},gobble=4]
    tex story
    dvips story
    pstopdf story.ps
  \end{lstlisting}

  \xxx
  \pause

  \begin{onlyenv}<presentation>
    \vskip -2.3cm
    \begin{tikzpicture}
      \draw[red,very thick] (0,0) -- (2.5cm, 2cm);
    \end{tikzpicture}
    \vskip .5cm
  \end{onlyenv}

  \begin{lstlisting}[language={},morekeywords={pdftex},gobble=4]
    pdftex story
  \end{lstlisting}
\end{Frame}

\begin{Frame}[t]{Ein \TeX-Dokument}{Quelltext \texttt{story.tex} und Dokument \texttt{story.pdf}}
  \lstinputlisting[firstline=7]{demo/story.tex}

  \begin{center}
    \includegraphics[width=10cm]{demo/story.pdf}
  \end{center}
\end{Frame}

\section*{Zusammenfassung}

\begin{frame}{Zusammenfassung}
  \begin{enumerate}
    \item Zusammenfassung.
  \end{enumerate}
\end{frame}